\chapter{Réduction du bruit}
\section{Séparer deux composantes}
\emph{Bruit de luminance} agit sur les fluctuations de clarté ;
\emph{Bruit de couleur} vise les variations chromatiques parasites.
Les deux contrôles vont de 0 à 100, référence zéro. Leur présence ne signifie
pas que le moteur utilise une IA équivalente à celle d’un autre logiciel.

\section{Méthode d’utilisation}
Corriger d’abord l’exposition utile et examiner une zone sombre relativement
uniforme, puis une zone texturée. Augmenter progressivement la réduction
de couleur si des taches colorées sont visibles. Ajuster ensuite la luminance
en vérifiant les cheveux, le tissu ou les détails fins. Revenir en vue ajustée
pour juger la sensation globale.

\paragraph{Exemple.}
Sur une photo à ISO élevés, ne pas conclure à partir du seul fond sombre.
Comparer également une texture et un contour net. Utiliser \touche{Ctrl + 1}
en gardant à l’esprit la résolution de l’aperçu, puis vérifier un export pleine
résolution pour une décision finale sur le détail.

\attention{Un débruitage plus fort peut supprimer de vrais détails.
Une référence Lightroom obtenue avec débruitage IA ne constitue pas une
comparaison à algorithmes identiques.}
\gestesreglages
